\section{Appendix 4.A: Discretization}\label{ch:discretization}
This appendix presents a general approach for discretizing a large class of continuum equilibrium problems that are governed by partial differential equations to formulate a representative discrete internal force vector $\intern$, as required by the \glspl{predictor-t1} of this chapter. The procedure generalizes the classical finite element method which is typically formulated on Sobolev vector spaces \citep{hughes1987finite} to the setting of nonlinear manifolds{\footnote{A manifold is a mathematical representation of a curved surface. Manifolds naturally arise as solution sets of constrained systems of equations and as graphs of functions.}}, which are encountered in the geometrically nonlinear analysis of plates, rods, and shells.
In all such problems, equilibrium can be stated in terms of nonlinear partial differential equations of the form:
\begin{equation}
\rho \ddot{\boldsymbol{u}} = \boldsymbol{f} -  \mathbb{B}_\chi\boldsymbol{\sigma},
\label{eq:strong}
\end{equation}
where \(\mathbb{B}_\chi\) is a nonlinear
differential operator and both the generalized stress \(\boldsymbol{\sigma}\) and the body force \(\boldsymbol{f}\) are \emph{continuous} functions defined over a body \(\mathcal{B}\) placed in space by a
map \(\chi\). In general \(\mathbb{B}\) may depend nonlinearly on \(\chi\),
but always acts linearly on \(\boldsymbol{\sigma}\). 
% This class of problems is formally summarized in Box 1.
% Problems with the form of
% \cref{eq:strong} describe the equilibrium of continuua, rods, and shells, whose
% configurations \(\chi\) may comprise a \emph{nonlinear} space
% \(\mathscr{C}\).
% Additional resources for this material are contained in
% \cite{marsden1994mathematical,zienkiewicz2014finite,ogden1997nonlinear} and \cite{shilov1996elementary}.

\hypertarget{sec:discretization}{%
\subsection{Approximation}\label{sec:discretization}}

The structure outlined in Box 1 implies a weak statement
of \cref{eq:strong} that is centered around the Galerkin \emph{residual functional}:
%
\begin{equation}
\mathcal{G}\left[\boldsymbol{\eta}\right] \triangleq \left< \boldsymbol{\eta}, \, \boldsymbol{f}\right>_\chi - \left< \boldsymbol{\eta}, \, \mathbb{B}_\chi \boldsymbol{\sigma}\right>_\chi.
\label{eq:galerkin}
\end{equation}
%
When specialized to classical continuum mechanics, this is equivalent to the classical virtual work functional (before integration by parts):
\begin{equation}
\int \delta\boldsymbol{u} \cdot \operatorname{div} \boldsymbol{\sigma}\, d\Omega +  \int \delta \boldsymbol{u} \cdot \boldsymbol{f}\, d\Omega, 
\label{eq:VW}
\end{equation}
where the virtual displacements \(\delta \boldsymbol{u}\) represent the test basis \(\boldsymbol{\eta}\) and the equilibrium operator $\mathbb{B}$ is represented by the (negative) divergence $-\operatorname{div}$. Note that the divergence always acts linearly on \(\boldsymbol{\sigma}\), but \(\boldsymbol{\sigma}\) may depend nonlinearly on the displacement \(\boldsymbol{u}\), resulting in a nonlinear problem.
The developments that follow
are kept applicable to a much wider class of mechanics models by retaining 
the inner product \(\left<\cdot,\,\cdot\right>\) and operator \(\mathbb{B}\) notations.
Note that in the following, 
the subscript \(\chi\) on \(\mathbb{B}_\chi\) and
\(\left<\cdot,\,\cdot\right>_\chi\) are dropped 
for brevity % {of notation}
and the body force \(\boldsymbol{f}\)
is not considered.
\iftrue
\begin{marsdenbox}[BOX 1\quad Continuous Mechanics]
\begin{description}
% \tightlist
\item[Body]
A body \(\mathcal{B}\) is a set of points \(\boldsymbol{\xi}\), and a placement
\(\mathcal{B}_\chi\) of \(\mathcal{B}\) is a Riemannian manifold
\(\mathcal{B}_\chi\) with boundary \(\partial \mathcal{B}\) and tangent
bundle \(T\mathcal{B}_\chi\) embedded in an ambient Euclidean space
\(\mathcal{E}\). The map
\(\chi\colon \mathcal{B}\rightarrow \mathcal{B}_\chi\) associates the
state of a material point \(\boldsymbol{\xi}\) to the placement
\(\mathcal{B}_\chi\). A reference configuration
\(\chi_0\colon \mathcal{B}\rightarrow \mathcal{B}_{\chi_0}\) is often
implied and the pull-back by the composition \(\chi\circ\chi_0^{-1}\) is
denoted by \(\boldsymbol{A}\).
% \tightlist
\item[Energy]
The set \(\mathscr{C}\) of configurations \(\chi\) forms a smooth
 manifold. A configuration \(\chi\) is identified by a couple
\((\boldsymbol{x}, \boldsymbol{\Lambda})\) where \(\boldsymbol{x}\in\mathscr{C}_x\) is a
vectorial configuration variable and
\(\boldsymbol{\Lambda}\in\mathscr{C}_{\scriptscriptstyle{\Lambda}}\) a
non-vectorial variable. These enter independently in a Riemannian metric
on the tangent space \(T_\chi\mathscr{C}\) to \(\mathscr{C}\) at
\(\chi\), which takes the assumed form: 
\begin{equation}
  \left< \cdot , \, \cdot\right>_{\chi} = \left< \cdot , \, \cdot\right>_{x} + \left< \cdot , \, \cdot\right>_{\scriptscriptstyle{\Lambda}} .
  \label{eq:conf-pair}
\end{equation}
%
\item[Stress]
The spatial stress \(\boldsymbol{\sigma}\) in \cref{eq:strong} is related to a
material stress \(\boldsymbol{s}\) which is defined with respect to inner products 
in \(\mathscr{C}\) at the current \(\chi\) and reference
\(\chi_0\) configurations:
\begin{equation}
  \left<\boldsymbol{c},\, \boldsymbol{\sigma}\right>_{\chi}
  =\left< \boldsymbol{A}\boldsymbol{c},\, \boldsymbol{s}\right>_{\chi_0}
  = \left<\boldsymbol{c},\, \boldsymbol{A}_*\boldsymbol{s}\right>_{\chi}
  %=\left< \boldsymbol{A}\boldsymbol{c},\, \boldsymbol{s}\right>_{\chi_0}\qquad\text{ for any }\boldsymbol{c},
  \label{eq:def-Astar}
\end{equation}
for any \(\boldsymbol{c}\).
% with \(\boldsymbol{A}_*\) denoting the Piola transformation. 
Similarly, a dual
\(\mathbb{B}^*\) to the equilibrium operator \(\mathbb{B}\) in
\cref{eq:strong} is defined with respect to \cref{eq:conf-pair}:
\begin{equation}
  \left<\boldsymbol{\eta},\, \mathbb{B}_\chi\boldsymbol{\sigma}\right>_\chi = \left<\mathbb{B}_\chi^*\boldsymbol{\eta},\, \boldsymbol{\sigma}\right>_\chi.
  \label{eq:def-adjoint}
\end{equation}
% \tightlist
\item[Strain]
The material strain \(\boldsymbol{e}\) is compatible with
\(\mathbb{B}^*_\chi\) through the differential relation: 
\begin{equation}
  D\boldsymbol{e}\cdot\boldsymbol{u} \equiv \boldsymbol{A}\mathbb{B}^*_\chi \boldsymbol{u}.
  \label{eq:vstrain}
\end{equation}
%
\item[Constitution]
The constitutive response is defined between the material stress
\(\boldsymbol{s}\) and the material strain \(\boldsymbol{e}\) with material \(\boldsymbol{C}_s\) and
spatial \(\boldsymbol{C}_\sigma\) moduli defined by:
\begin{equation}
  D \boldsymbol{s}\cdot \boldsymbol{u} = \boldsymbol{C}_s \, D \boldsymbol{e}\cdot {\boldsymbol{u}}
  \qquad\text{ and }\qquad
  \boldsymbol{C}_\sigma \triangleq \boldsymbol{A}_*\boldsymbol{C}_s\boldsymbol{A}.
  \label{eq:def-moduli}
\end{equation}
\end{description}
\end{marsdenbox}
\fi

In nonlinear finite element analysis,
the solution algorithm searches for a configuration \(\chi\) in the space \(\mathscr{C}\) that satisfies the
equilibrium relations in
\cref{eq:strong} by applying a sequence of small
perturbations \(\{\boldsymbol{u}^{(i)}\}\) to an initial guess \(\chi^{(0)}\).
Each perturbation \(\boldsymbol{u}^{(i)}\) is determined by the linearization of
\cref{eq:galerkin} at %around
the last trial configuration \(\chi^{(i)}\) of the iterative process.
At each iteration, the tangent space \(T_{\chi^{(i)}}\mathscr{C}\) contains infinitely many modes
for representing \(\boldsymbol{u}^{(i)}\) in %for
updating the configuration \(\chi^{(i)}\).
%which might be represented in.
% These modes represent \PFF{the} available directions for the
% next trial configuration in the space \(\mathscr{C}\).
This appendix develops the general equations for % this procedure
the procedure of updating the trial configuration \(\chi^{(i)}\)
by first linearizing \cref{eq:galerkin} in Section~\ref{sec:linearization} and then developing a
finite-dimensional representation of the linearized residual in
Section~\ref{sec:galerkin-discretization}.

\hypertarget{sec:linearization}{%
\subsection{Linearization}\label{sec:linearization}}

The linearization of the residual functional \(\mathcal{G}\)
in \cref{eq:galerkin} is akin to a first-order Taylor series expansion
given by:
%
\begin{equation}
\mathscr{L} \mathcal{G}\left[\boldsymbol{\eta} \right]
=\mathcal{G}\left[\boldsymbol{\eta}\right] + \nabla_{\boldsymbol{u}} \mathcal{G}\left[\boldsymbol{\eta}\right],
\label{eq:galerkin-linearized}
\end{equation}
where
\(\nabla_{\boldsymbol{u}}\mathcal{G}[\boldsymbol{\eta}]\) is the covariant derivative of
the functional \(\mathcal{G}\).
This is given by the formula for the covariant differentiation of co-vectors: % \citep{romano2005tangentb, simo1992symmetric}:
\begin{equation}
\nabla_{\boldsymbol{u}}\mathcal{G}\left[\boldsymbol{\eta}\right] =
D_{}\mathcal{G}\left[\boldsymbol{\eta}\right]\cdot \boldsymbol{u}-\mathcal{G}\left[\nabla_{\boldsymbol{u}} \boldsymbol{\eta}\right],
\label{eq:galerkin-linearized1}
\end{equation}
where \(D\mathcal{G}\left[\boldsymbol{\eta}\right]\cdot\boldsymbol{u}\) is the
Gateaux differential of the residual component
\(\mathcal{G}\left[\boldsymbol{\eta}\right]\) in the direction
\(\boldsymbol{u} \in T_\chi \mathscr{C}\) and
\(\mathcal{G}\left[\nabla_{\boldsymbol{u}}\boldsymbol{\eta}\right]\) accounts for the
nonlinearity of the manifold.
The covariant term \(\mathcal{G}\left[\nabla_{\boldsymbol{u}}\boldsymbol{\eta}\right]\) is commonly
neglected in the literature, which amounts to selecting an implicit connection, and often produces an asymmetric tangent stiffness matrix.

At the \(i^{\text{ th}}\) iteration, the differential is
defined in terms of an integral curve
\(\tau \mapsto {\chi}_\tau\) through \(\mathscr{C}\) with
the properties:
%
\begin{equation}
\chi_{\tau}(0) = \chi^{(i)} \in \mathscr{C}
\qquad\text{ and }\qquad
  \dot{\chi}_\tau(0) %\biggr|_{\tau=0} 
  = \boldsymbol{u}^{(i)} \in T_\chi\mathscr{C}.
\label{eq:chi-curve}
\end{equation} 
To produce distinct material and geometric tangents,
\cref{eq:galerkin} is 
separated into terms with %first represented as a product of
purely material and geometric quantities.
This is achieved by % successively
transposing the operators \(\mathbb{B}\) and \(\boldsymbol{A}_*\)
with the help of %by
\cref{eq:def-Astar} and  \cref{eq:def-adjoint}
to shift the %, which shifts these
non-material terms %quantities
in \cref{eq:galerkin} to the first argument of the inner product, leaving only the material stress \(\boldsymbol{s}\) as the second argument:
%
\begin{equation}
\begin{aligned}
\mathcal{G} [\boldsymbol{\eta}]
&= \left< \boldsymbol{\eta},\,  \mathbb{B} \boldsymbol{\sigma}\right>  = \left< \mathbb{B}^*_{} \boldsymbol{\eta}, \, \boldsymbol{\sigma}\right> = \left< \boldsymbol{A}\mathbb{B}^*_{} \boldsymbol{\eta}, \, \boldsymbol{s}\right>.
\end{aligned}
\label{eq:galerkin-material}
\end{equation}
Linearization now proceeds in the standard fashion with the application of the product rule to the last expression of \cref{eq:galerkin-material}:

\begin{equation}
\begin{array}{rccc}
D\mathcal{G}[\boldsymbol{\eta}]\cdot \boldsymbol{u}
=& \frac{\mathrm{d}}{\mathrm{d}\tau} \mathcal{G}_\tau [\boldsymbol{\eta}] ~\biggr|_{\tau=0} \\[0.3cm]
  % =& \frac{\mathrm{d}}{\mathrm{d}\eta} \left.\left< \boldsymbol{A}^\mathrm{t}\mathbb{B}^*_{} \boldsymbol{\eta} , \, \boldsymbol{s}\right> \right|_{\eta=0} \\[0.3cm]
  =&  \left< \boldsymbol{A}\mathbb{B}^*_{} \boldsymbol{\eta} , \, D\boldsymbol{s}\cdot \boldsymbol{u}\right> &+& \left< D_{}[ \boldsymbol{A} \mathbb{B}^*_{}\boldsymbol{\eta}]\cdot \boldsymbol{u}, \, \boldsymbol{s}\right>. \\[0.3cm]
  \triangleq& \mathcal{K}_{M}\left[\boldsymbol{\eta}, \boldsymbol{u}\right] &+& \mathcal{K}_{G}\left[\boldsymbol{\eta}, \boldsymbol{u}\right], \\
  % \triangleq& &\mathcal{K}\left[\boldsymbol{\eta}, \boldsymbol{u}\right]
\end{array}
\label{eq:def-K}
\end{equation} 
where \(\mathcal{K}_{M}\) and \(\mathcal{K}_{G}\) are referred to as the \emph{material} and \emph{geometric} stiffness \emph{operators} of the problem. 
The stress derivative \(D\boldsymbol{s}\cdot \boldsymbol{u} \) that appears 
in \(\mathcal{K}_{M}\) above 
is taken fixed in the reference configuration
and expanded by the chain rule to give:
\begin{equation}
\begin{aligned}
D\boldsymbol{s}(\boldsymbol{e}) \cdot \boldsymbol{u} &= \left.\frac{\mathrm{d}}{\mathrm{d} \tau} \left. \boldsymbol{s}\circ \boldsymbol{e} \right. ~\right|_{\tau=0} \\
&= \boldsymbol{C}_s \, D \boldsymbol{e} \cdot \boldsymbol{u} \\
&= \boldsymbol{C}_s \, \boldsymbol{A}\mathbb{B}^*_{} \boldsymbol{u},
\end{aligned}
\end{equation}
where \cref{eq:def-moduli}
supplies %has supplied
the modulus \(\boldsymbol{C}_s\) and
\cref{eq:vstrain} the strain linearization
\(\boldsymbol{A}\mathbb{B}^*_{} \boldsymbol{u}\).
Substitution back into \cref{eq:def-K} furnishes the final expression for the material stiffness
\(\mathcal{K}_M\) operator:

\begin{equation}
\begin{aligned}
\mathcal{K}_{M}\left[\boldsymbol{\eta}, \boldsymbol{u} \right]
%&= \left< \boldsymbol{A}^\mathrm{t}\mathbb{B}^*_{} \boldsymbol{\eta} , \, D\boldsymbol{s}\cdot \boldsymbol{u}\right> \\
&= \left< \boldsymbol{A}\mathbb{B}^*_{} \boldsymbol{\eta} , \, \boldsymbol{C}_s \, \boldsymbol{A}\mathbb{B}^*_{} \boldsymbol{u}\right> %\\
= \left<\mathbb{B}^*_{} \boldsymbol{\eta} , \, \boldsymbol{C}_\sigma \, \mathbb{B}^*_{} \boldsymbol{u}\right>, \\
\end{aligned}
\label{eq:Km}
\end{equation}
where \cref{eq:def-moduli} is used to replace \(\boldsymbol{C}_s\) by \(\boldsymbol{C}_{\sigma}\)
after transferring \(\boldsymbol{A}\) in the first inner product argument to the second argument
with \cref{eq:def-Astar}.

\iftrue
The geometric stiffness \(\mathcal{K}_{G}\) is characterized by the directional derivative
\(D[\boldsymbol{A}\mathbb{B}^*\boldsymbol{\eta}]\cdot \boldsymbol{u}\). Application of the
product rule furnishes:

\begin{equation}
\begin{aligned}
\mathcal{K}_{G}\left[\boldsymbol{\eta},\boldsymbol{u}\right]
&= \left<\partial_{\boldsymbol{u}}\left[\boldsymbol{A}\right]\mathbb{B}^*\boldsymbol{\eta} ,\, \boldsymbol{s}\right> +  \left<\boldsymbol{A}\partial_{\boldsymbol{u}}\left[\mathbb{B}^*\boldsymbol{\eta}\right] ,\, \boldsymbol{s}\right>.
\end{aligned}
\label{eq:def-Kg}
\end{equation}
The first term on the right arises from
the push-forward between the material and the spatial representation while the
second term captures the dependence of the operator \(\mathbb{B}^*\) on the configuration \(\chi\).
In a small deformation setting where the
material and spatial representations are assumed to coincide,
both terms will naturally vanish. %both of these terms naturally vanish.
\fi
\hypertarget{sec:galerkin-discretization}{%
\subsection{Discretization}\label{sec:galerkin-discretization}}

Because \cref{eq:galerkin-linearized} and \cref{eq:def-K} are linear in both
\(\boldsymbol{\eta}\) and \(\boldsymbol{u}\), a discrete algebraic problem
%follows %immediately for
arises directly when the discretized vectors \(\boldsymbol{\eta}^h\) and \(\boldsymbol{u}^h\)
%that are 
are represented as a linear combination of \emph{simple} basis functions
that only depend
on the coordinate \(\boldsymbol{\xi}\).
%This is the case when the basis functions only depend
However, for problems formulated on manifolds,
a simple interpolation is often inappropriate.
Instead it is often desirable to apply a \emph{configuration-dependent} interpolation, where
\(\boldsymbol{\eta}^h\) and \(\boldsymbol{u}^h\)
depend nonlinearly on the nodal values 
\citep{romero2004interpolation, cardona1988beam, simo1990stress}.
In this case, discrete expressions for \cref{eq:galerkin-linearized} and
\cref{eq:def-K} follow by constructing
the curve \(\chi_\tau\) in \cref{eq:chi-curve}
in terms of discrete nodal parameters \(\{\boldsymbol{u}_b\}\)
furnishing the discrete tangent space \(T_\chi\mathscr{C}^h\) 
where elements take the form:
\begin{equation*}
%\begin{aligned}
%\left.\frac{\mathrm{d}}{\mathrm{d}\tau} \chi_\tau \right|_{\tau=0} &= \sum_a \partial_a\boldsymbol{\eta}^h \cdot \boldsymbol{\eta}_a \\
%&\triangleq \boldsymbol{\eta}^\ell
%\end{aligned}
%\qquad\text{ and }\qquad
\begin{aligned}
\boldsymbol{u}^\ell \triangleq \left.\frac{\mathrm{d}}{\mathrm{d}\tau} \chi_\tau \right|_{\tau=0} 
  &= \sum_b \left.\frac{\mathrm{d}}{\mathrm{d}\tau}\boldsymbol{u}^h(\tau \boldsymbol{u}_b)\right|_{\tau=0}. \\
\end{aligned}
\end{equation*}
This motivates the definition of \emph{effective interpolation bases}
\(\left\{\boldsymbol{\varphi}^{}_{\eta a}\right\}\) and
\(\left\{\boldsymbol{\varphi}^{}_{u b}\right\}\):
\begin{equation}
  \boldsymbol{\varphi}^{}_{\eta a}(\boldsymbol{\xi}) \triangleq \left.\frac{\mathrm{d}}{\mathrm{d}\tau}\boldsymbol{\eta}^h (\boldsymbol{\xi}; \tau\boldsymbol{u}_a)\right|_{\tau=0}
\qquad\text{ and }\qquad
  \boldsymbol{\varphi}^{}_{u b}(\boldsymbol{\xi}) \triangleq \left.\frac{\mathrm{d}}{\mathrm{d}\tau}\boldsymbol{u}^h(\boldsymbol{\xi}; \tau  \boldsymbol{u}^{}_b)\right|_{\tau=0}.
\label{eq:interp}
\end{equation}

\hypertarget{resisting-forces}{%
\subsubsection{Resisting forces}\label{resisting-forces}}

%The coordinate expression of on the basis
The evaluation of \cref{eq:galerkin}
on the linearized variation \(\boldsymbol{\eta}^\ell\)
furnishes a representation on the basis 
\(\boldsymbol{\varphi}_{\eta a}\) which
follows 
from the linearity of the inner product:
\begin{equation}
  \begin{aligned}
  \mathcal{G}\left[\boldsymbol{\eta}^\ell\right] &= \left< \boldsymbol{\eta}^h ,\, \mathbb{B}\boldsymbol{\sigma} \right> %\\ & 
  = \sum_a \left< \boldsymbol{\varphi}^{}_{\eta a} \boldsymbol{\eta}_a ,\, \mathbb{B}\boldsymbol{\sigma} \right>. \\
  \end{aligned}
\end{equation}
Integration by parts is once again expressed by shifting
from the operator \(\mathbb{B}\) in the second argument to
\(\mathbb{B}^*\) in the first:
%integrating by parts furnishes:
\begin{equation}
  \begin{aligned}
  \mathcal{G}\left[\boldsymbol{\eta}^\ell\right]&= \sum_a \left<\mathbb{B}^*_{} \left[\boldsymbol{\varphi}^{}_{\eta a} \boldsymbol{\eta}_a\right] ,\, \boldsymbol{\sigma} \right>   \\
  %&= \sum_a \left<\mathbb{B}^*_{} \left[\boldsymbol{\varphi}^{}_{\eta a} \boldsymbol{\eta}_a\right] ,\, \boldsymbol{\sigma} \right>  + \mathcal{G}_{\partial} \\
  &= \sum_a \left< \boldsymbol{B}_a^{\mathrm{t}}\boldsymbol{\eta}_a ,\,   \boldsymbol{\sigma} \right>  \\
  &= \sum_a \left< \boldsymbol{\eta}_a ,\,  \boldsymbol{B}_a \boldsymbol{\sigma} \right>, \\
  %&= \sum_a \left< \boldsymbol{\eta}_a ,\, \left[\mathbb{B}^*_{} \boldsymbol{\varphi}^{}_{\eta a}\right]^{\mathrm{t}}  \boldsymbol{\sigma} \right> + \mathcal{G}_{\partial}. \\
  \end{aligned}
  \quad\text{ with }\quad
  \boldsymbol{B}_a^{\mathrm{t}} \triangleq \mathbb{B}^*\boldsymbol{\varphi}^{}_{\eta a}
\label{eq:galerkin-ibp}
\end{equation}
where \(\boldsymbol{B}_a^{\mathrm{t}}\) is
the linear operator that results
from the application of 
\(\mathbb{B}^*\) to \(\boldsymbol{\varphi}^{}_{\eta a}\).
Finally, recall that for any orthonormal
basis \(\{\mathbf{e}_k\}\) spanning the degrees of freedom in
\(T\mathcal{B}_{\chi}\), an identity
is furnished by 
the summation \(\mathbf{e}_k \otimes \mathbf{e}_k\).
This is used in \cref{eq:galerkin-ibp} to give:

\begin{equation}
\begin{aligned}
\mathcal{G}\left[\boldsymbol{\eta}^\ell\right]
&= \sum_{a,k} \left<\left(\mathbf{e}_k\otimes\mathbf{e}_k\right) \boldsymbol{\eta}_a,\, \boldsymbol{B}_a  \boldsymbol{\sigma} \right> \\
&= \sum_{a,k} \boldsymbol{\eta}_a \cdot \mathbf{e}_k  \left<\mathbf{e}_k ,\, \boldsymbol{B}_a  \boldsymbol{\sigma} \right>, \\
%&= \sum_a \boldsymbol{\eta}_a\cdot \int \left[\mathbb{B}^*_{} \phi_a \boldsymbol{1}\right]^{\mathrm{t}} \boldsymbol{\sigma} \, d\boldsymbol{\xi} \\
%&= \sum_a \boldsymbol{\eta}_a\cdot \intern_a
\end{aligned}
\label{eq:global-work-dot}
\end{equation}
where the summation \(\left<\mathbf{e}_k ,\,\boldsymbol{B}_a  \boldsymbol{\sigma} \right> \mathbf{e}_k\) 
represents the residual element nodal forces \(\{\intern_a\}\).
% \begin{equation}
% \intern_a \triangleq  \left<\mathbf{e}_k ,\,\boldsymbol{B}_a  \boldsymbol{\sigma} \right> \mathbf{e}_k .
% \label{eq:def-pa}
% \end{equation}

\hypertarget{tangent-stiffness}{%
\subsubsection{Tangent Stiffness}\label{tangent-stiffness}}

The material tangent \(\mathcal{K}_M\left[\boldsymbol{\eta}^\ell, \boldsymbol{u}^\ell\right]\)
depends on \(\boldsymbol{u}^\ell\) %under
{via} the operator \(\mathbb{B}^*\).
{On account of linearity}, % which by linearity
{the operation has} %furnishes 
the following representation:
\begin{equation}
\mathbb{B}^*_{} \boldsymbol{u}^\ell = \sum_b \boldsymbol{B}^\star_{b} {\boldsymbol{u}}_b
\qquad\text{ with }\qquad
\boldsymbol{B}^\star_b \triangleq \mathbb{B}^*\boldsymbol{\varphi}_{u b}.
\label{eq:def-Bb}
\end{equation} 
Similarly, %define
a discrete geometric operator \(\boldsymbol{G}_{ab}[\boldsymbol{s}]\)
is defined
as the partial derivatives of the product \(\boldsymbol{B}_a\boldsymbol{A}_*\boldsymbol{s}\)
with respect to the nodal parameters \(\boldsymbol{u}_b\),
with \(\boldsymbol{s}\) held constant:
%
\begin{equation}
\boldsymbol{G}_{ab}[\boldsymbol{s}] = \partial_b[\boldsymbol{B}_a\boldsymbol{A}_*\mid\boldsymbol{s}],
%\boldsymbol{G}_{ab}[\mathbf{e}_i]\triangleq \partial_b \left[\boldsymbol{A}\mathbb{B}^*\boldsymbol{\varphi}_{\eta a}\mathbf{e}_i\right]\boldsymbol{\varphi}_{u b}.
\label{eq:def-Gab}
\end{equation}
where the vertical bar indicates that the terms that follow are held constant under differentiation.
%
In terms of \(\boldsymbol{B}_a\)%(\cref{eq:galerkin-ibp})
,     \(\boldsymbol{B}^\star_b\)%(\cref{eq:def-Bb})
, and \(\boldsymbol{G}_{ab}\),   %(\cref{eq:def-Gab}) 
the Galerkin 
differential \(D\mathcal{G}\left[\boldsymbol{u}^\ell\right]\cdot\boldsymbol{\eta}^\ell\) of Equation (\ref{eq:def-K}) then has the representation:
\begin{equation} 
\begin{aligned}
\mathcal{K}_M\left[\boldsymbol{\eta}^\ell, \boldsymbol{u}^\ell\right] + \mathcal{K}_G\left[\boldsymbol{\eta}^\ell, \boldsymbol{u}^\ell\right] = \sum_{a,b} \boldsymbol{\eta}_a\cdot\mathbf{K}_{ab} \, \boldsymbol{u}_b,
\end{aligned}
\end{equation}
where the nodal stiffness matrices \(\{\mathbf{K}_{a b}\}\)
are furnished by summation on \(i\) and \(j\):
\begin{equation}
\begin{aligned}
\mathbf{K}_{ab} 
  &\triangleq 
  \left<\mathbf{e}_i ,\,\boldsymbol{B}^{}_a \boldsymbol{C}_{\sigma}\boldsymbol{B}^{\star}_b \mathbf{e}_j\right> \mathbf{e}_i\otimes\mathbf{e}_j
  +\left<\mathbf{e}_i ,\, \boldsymbol{G}_{ab}[\boldsymbol{s}]\mathbf{e}_j\right> \mathbf{e}_i\otimes\mathbf{e}_j .
%&\triangleq %\mathcal{K}_M\left[\boldsymbol{\varphi}_{\eta a}\mathbf{e}_k, \boldsymbol{\varphi}_{u b}\mathbf{e}_\ell\right]\\
%\left<\mathbf{e}_i ,\,\boldsymbol{B}^{}_a \boldsymbol{C}^{}_{\sigma}\boldsymbol{B}^{\star}_b \mathbf{e}_j\right> \mathbf{e}_i\otimes\mathbf{e}_j
%+\left<\boldsymbol{G}_{ab}[\mathbf{e}_i]\mathbf{e}_j
%,\, \boldsymbol{s}\right> \mathbf{e}_i\otimes\mathbf{e}_j .
\end{aligned}
\label{eq:def-kab}
\end{equation}


\hypertarget{remarks03}{%
\subsection{Summary}\label{remarks03}}

  The results of this appendix are collected in Box 2.
  The operators \(\boldsymbol{B}_a\),     \(\boldsymbol{B}^\star_b\), and \(\boldsymbol{G}_{ab}\) 
  in Box 2 are the essential components that define a finite element and furnish the element internal forces \(\intern\) and stiffness \(\mathbf{K}\).
  If \(\boldsymbol{\varphi}_{\eta a}\) and \(\boldsymbol{\varphi}_{u b}\) are
  equivalent, the discrete equilibrium operator \(\boldsymbol{B}_a\) 
  will be the transpose of the
  discrete strain operator \(\boldsymbol{B}^\star_a\)
  and the discrete material stiffness will be 
  symmetric.
% \item

  In a \emph{Bubnov-Galerkin} formulation, the iterative solution procedure
  converges to a configuration where the residual of \cref{eq:strong}
  is orthogonal to the chosen modes of perturbation, which may be represented in
  \(\boldsymbol{u}\).
  This variant is
  commonly referred to simply as \emph{the Galerkin} method and shares
  an intimate link with problems that can be derived as energy
  minimizers.
% \item

  Alternatively, a \emph{Petrov-Galerkin} formulation arises when
  variations \(\boldsymbol{u}^h\) and \(\boldsymbol{\eta}^h\) are permitted distinct
  representations. This can become preferable when a complicated basis
  is chosen for \(\boldsymbol{u}\), as it can greatly simplify the derivation of
  a consistent tangent stiffness.
% \end{itemize}

\begin{marsdenbox}[BOX 2\quad Finite Element Operators]
{\noindent

\begin{itemize}
\item
  A Galerkin approximation of %to
  \cref{eq:galerkin-linearized} is defined by effective interpolation bases 
  \(\left\{\boldsymbol{\varphi}^{}_{\eta a}\right\}\) and \(\left\{\boldsymbol{\varphi}^{}_{u b}\right\}\) where:
  \begin{equation*}
    \boldsymbol{\eta}^\ell(\boldsymbol{\xi}) = \sum_a \boldsymbol{\varphi}^{}_{\eta a}(\boldsymbol{\xi})\boldsymbol{\eta}^{}_a
    \qquad\text{ and }\qquad
    \boldsymbol{u}^\ell(\boldsymbol{\xi}) = \sum_b \boldsymbol{\varphi}^{}_{u b}(\boldsymbol{\xi}) \boldsymbol{u}^{}_b.
  \end{equation*}

\item
  Application to the linearized residual of \cref{eq:galerkin-linearized} furnishes \emph{three} characteristic operators:
  \begin{equation*}
  \begin{array}{lrll}
    %&\textbf{Symbol} &\textbf{Expression} & \textbf{Description} \\[1ex] %\hline
    \text{(i)}  &\boldsymbol{B}_a^{\mathrm{t}}     & = \mathbb{B}^*\boldsymbol{\varphi}^{}_{\eta a} & \text{Discrete equilibrium operator} \\[2ex]
    \text{(ii)} &\boldsymbol{B}^{\star}_b          & = \mathbb{B}^*\boldsymbol{\varphi}^{}_{u b}         & \text{Discrete strain operator} \\[2ex]
    \text{(iii)}&\boldsymbol{G}_{ab}[\boldsymbol{s}] & = \partial_b\left[\boldsymbol{B}_a\boldsymbol{A}_*\mid\boldsymbol{s}\right]%\boldsymbol{\varphi}_{u b}
    %\partial_b \left[\boldsymbol{A}\mathbb{B}^*\boldsymbol{\varphi}_{\eta a}\mathbf{e}_i\right]\boldsymbol{\varphi}_{u b} & 
    & \text{Discrete geometric operator} \\[2ex]
  \end{array}
  \end{equation*}

\iffalse
\begin{center}
\begin{tabular*}{250pt}{@{\extracolsep\fill}lll@{\extracolsep\fill}}%
\toprule
 & \textbf{Expression} & \textbf{Description} \\
\midrule
1. &  \(\boldsymbol{B}_a^{\mathrm{t}} \triangleq \mathbb{B}^*\boldsymbol{\varphi}^{}_{\eta a}\) & Discrete equilibrium operator \\
2. &  \(\boldsymbol{B}^{\star}_b \triangleq \mathbb{B}^*\boldsymbol{\varphi}^{}_{u b}\)       & Discrete strain operator \\
3. &  \(\boldsymbol{G}_{ab}[\mathbf{e}_i] \triangleq \partial_b \left[\boldsymbol{A}\mathbb{B}^*\boldsymbol{\varphi}_{\eta a}\mathbf{e}_i\right]\boldsymbol{\varphi}_{u b}\) & Discrete geometric tangent operator\\
%4. &  \(\boldsymbol{H}_{ab} \triangleq \boldsymbol{A}D\left[\mathbb{B}^*[\boldsymbol{\varphi}_{ub}\mathbf{e}_{k}]\right]\cdot\boldsymbol{\varphi}^{}_{u b}\mathbf{e}_{\ell}\) & Second geometric tangent \\
\bottomrule
\end{tabular*}
\end{center}
\fi

\item 
  These operators determine the \emph{nodal residual forces} and the \emph{tangent stiffness}, which  in terms of an orthonormal basis \(\left\{\mathbf{e}_k\right\}\) spanning the degrees of freedom in \(T\mathcal{B}_{\chi}\) are furnished by:
  \begin{equation}
    \begin{aligned}
    \intern_a =& \left<\mathbf{e}_i ,\, \boldsymbol{B}_a\boldsymbol{\sigma}\right>\mathbf{e}_i, \\
    \extern_a =& \left<\boldsymbol{\varphi}_{\eta a} \mathbf{e}_k,\, \boldsymbol{f}\right>\mathbf{e}_k,
    \end{aligned}
    \label{eq:def-pa}
  \end{equation}
  and
  \begin{equation*}
    \begin{aligned}
    \mathbf{K}_{ab} 
     =& \left<\mathbf{e}_i ,\,\boldsymbol{B}^{}_a \boldsymbol{C}_{\sigma}\boldsymbol{B}^{\star}_b \mathbf{e}_j\right> \mathbf{e}_i\otimes\mathbf{e}_j
      +\left<\mathbf{e}_i ,\, \boldsymbol{G}_{ab}[\boldsymbol{s}]\mathbf{e}_j\right> \mathbf{e}_i\otimes\mathbf{e}_j. 
    \end{aligned}
  \end{equation*}
\end{itemize}
}
\end{marsdenbox} 

